\documentclass[10pt,a4paper]{amsart}
\usepackage[margin=19mm]{geometry}
\usepackage{amsmath,amssymb,mathtools}
\usepackage[scr=rsfs,cal=euler]{mathalfa}
\usepackage{xfrac}
\usepackage[unicode,hidelinks,bookmarks=false]{hyperref}
\newcommand{\ie}{i.e.\ }
\newcommand{\cf}{cf.\ }
\newcommand{\resp}{respectively\ }
\newcommand{\Subsection}{\subsection}
\providecommand{\nobreakdash}{\nobreak\hbox{-}}
\setlength{\emergencystretch}{2em}
\multlinegap0pt
\usepackage{bm}
\usepackage{bbm}
\usepackage{dsfont}
\usepackage{xspace}
\usepackage{cancel}
\usepackage{mathtools}
\usepackage{amsthm}
\makeatletter
\@ifundefined{theorem}{\newtheorem{theorem}{Theorem}}{}
\@ifundefined{lemma}{\newtheorem{lemma}[theorem]{Lemma}}{}
\makeatother
\title[Nearly optimal spectral gaps for random Belyi surfaces]{Nearly optimal spectral gaps for random Belyi surfaces}
\author{Yang Shen, Yunhui Wu}
\date{Source: arXiv:2511.02517v3 (CC BY 4.0)}
\begin{document}
\maketitle

\noindent\emph{Source and licence.} Nearly optimal spectral gaps for random Belyi surfaces. Version arXiv:2511.02517v3 (CC BY 4.0), distributed under the Creative Commons Attribution 4.0 International licence, as recorded in the arXiv licence metadata for this exact version and shown on the abstract page https://arxiv.org/abs/2511.02517v3. Full rights notice: licenses/arxiv-2511.02517-CC-BY-4.0.txt.\par\smallskip

\noindent\textit{Editorial scope.} Counted unit: the Lemma whose source label is \texttt{l-decom} in the article (source0.tex lines 1334--1349, source label \texttt{l-decom}), delivered together with the article's own complete original proof of it (source0.tex lines 1350--1421, one \texttt{proof} environment of 72 lines). No mathematics is written, repaired or re-derived: statement and proof are the article's own text line for line. Required context retained uncounted: none. Every definition, display and label of those lines is retained unchanged. Omitted from this package and not redistributed: 1--1333, 1422--1769. Nothing inside a retained excerpt is omitted or altered: every retained body is delivered exactly as the article's lines are written, so no omission receipt of any kind accompanies this package. Citations: the retained excerpts carry no citation command at all, so no bibliography entry of the article is transcribed and no reference is redistributed. Reference adaptation: none was needed and none was made; every reference inside the delivered unit resolves inside it, so no unresolved reference or citation is produced. Printed slips kept literal: no printed word, symbol, hypothesis or number of the counted statement or of its proof was altered. Layout and class: the article's own class is \texttt{amsart} with packages including amssymb, euscript, tikz, units, hyperref, verbatim, booktabs, array, boldline, color, cleveref, colonequals, tikz-cd, appendix; the wrapper is the lane's pinned \texttt{amsart} editorial wrapper at 10pt on A4, which restates the theorem-like environments the retained text needs and adds the article's own macro definitions that the retained text needs (none), with extra wrapper packages (bm, bbm, dsfont, xspace, cancel, mathtools). The delivered numbering is the unit's own. Assets: no retained span contains a figure environment, a graphics-inclusion command, a PDF-page inclusion, a table or a TikZ picture; no image file is copied into this package, the package declares no asset dependency and no image file is redistributed with it. Licence: version arXiv:2511.02517v3 of this article is distributed under the Creative Commons Attribution 4.0 International licence, as recorded in the arXiv licence metadata for this exact version and shown on the abstract page https://arxiv.org/abs/2511.02517v3. A full rights notice is delivered at licenses/arxiv-2511.02517-CC-BY-4.0.txt. Characters: every retained excerpt is pure ASCII, so the delivered engine, XeLaTeX with its default Latin Modern text font, reports no missing character.
    \begin{lemma}\label{l-decom}
        Fix a finite subset $\mathcal{S}\subset\textnormal{PSL}(2,\mathbb{Z})$ such that $\max\limits_{\gamma\in\mathcal{S}}|\gamma|\leq d$, and let $\omega$ be any not-necessary-reduced word of length $p$ with letters in $\mathcal{S}$, 
        \begin{enumerate}
            \item if $\omega$ is conjugate to one of $\{b,c,c^2\}$ as an element in $\textnormal{PSL}(2,\mathbb{Z})$, then there exist $1\leq t_1\leq t_2\leq p$,  elements $\xi,v_1,v_2\in\textnormal{PSL}(2,\mathbb{Z})$, and $h\in\{b,c,c^2\}$ with $|v_i|\leq d$ $(i=1,2)$ such that
            $$\omega_{[1,t_1)}v_1\equiv \xi,\ v_1^{-1}\omega_{[t_1,t_2)}v_2\equiv h,\ v_2^{-1}\omega_{[t_2,p]}\equiv \xi^{-1}.$$
            \item If $\omega$ is conjugate to $h_0^m\ (m\geq 2)$ as an element in $\textnormal{PSL}(2,\mathbb{Z})$ for some $h_0\in\textnormal{PSL}(2,\mathbb{Z})$, where $h_0$ is primitive and 
            $$h_0=bc^{i_1}\cdots bc^{i_k}$$
            for some $k\geq 1$ and $i_j\in\{1,2\}\ (1\leq j\leq k)$,
            then there exist $1\leq t_1\leq t_2\leq t_3\leq t_4\leq p$ and elements $\xi,h,v_1,...,v_4\in\textnormal{PSL}(2,\mathbb{Z})$ with $|v_i|\leq d+1\ (1\leq i\leq 4)$ such that
            $$\begin{matrix}
             \omega_{[1,t_1)}v_1\equiv \xi,& v_1^{-1}\omega_{[t_1,t_2)}v_2\equiv v_2^{-1}\omega_{[t_2,t_3)}v_3\equiv h,\\
                v_3^{-1}\omega_{[t_3,t_4)}v_4\equiv h^{m-2},& v_4^{-1}\omega_{[t_4,p]}\equiv \xi^{-1}.
            \end{matrix}$$
        \end{enumerate}
        Here, $``\equiv"$ means that two words are the same element in $\textnormal{PSL}(2,\mathbb{Z})$.
    \end{lemma}

    \begin{proof}
        Firstly by our assumption one may write
        \[\omega=\omega_1\omega_2\cdots\omega_p,\]
        where each $\omega_i=e_i^1\cdots e_i^{l_i}\in\mathcal{S}$ with $1\leq l_i\leq d$ and $e_i^j \in \{b,c\}$. So we write 
        $$\omega_1\omega_2\cdots\omega_p=\omega=e_1e_2\cdots e_l,$$
        where $l=\sum_{i=1}^{p}l_i\leq pd$ and $e_i\in\{b,c\}\ (1\leq i\leq l)$. For any non-unitary element $\theta\in\textnormal{PSL}(2,\mathbb{Z})$, it could be represented as 
        $$\theta=f_1\cdots f_t,$$
        where $f_i\in\{b,c,c^2\}$ and $\{f_i,f_{i+1}\}=\{b,c\}$ or $\{b,c^2\}$ for each $1\leq i\leq t-1$. We call it a $(b,c)-$representation of $\theta$.  In remaining part of the proof, for $f=b,\ c$ and $c^2$, we set $f^{-1}=b,\ c^2$ and $c$ respectively.
        
       (1). We first consider the case that $\omega$ is conjugate to $b$. Then there exists $\xi_0\in\textnormal{PSL}(2,\mathbb{Z})$  such that $\omega\equiv \xi_0 b \xi_0^{-1}$. Assume 
       $$\xi_0=f_1\cdots f_t$$
       is a $(b,c)-$representation of $\xi_0$. 
       If $f_t=b$, then $f_{t-1}=c$ or $c^2$ and  
       $$\omega\equiv f_1\cdots f_tbf_t^{-1}\cdots f_1^{-1}\equiv f_1\cdots f_{t-1}b f_{t-1}^{-1}\cdots f_1^{-1},$$
       where the word in the right hand side is reduced. Take $\xi=f_1\cdots f_{t-1}$, then $\omega\equiv \xi b \xi^{-1}$, and there exist $1\leq j_1\leq j_2\leq l$ such that 
        $$e_1e_2\cdots e_{j_1}\equiv \xi\text{ and }e_1e_2\cdots e_{j_2}\equiv \xi b.$$
        Hence 
        $$e_{j_1+1}\cdots e_{j_2}\equiv b\text{ and }e_{j_2+1}\cdots e_l\equiv \xi^{-1}.$$
        Assume $e_{j_1}$ and $e_{j_2}$ are contained in $\omega_{t_1}$ and $\omega_{t_2}$ $(1\leq t_1\leq t_2\leq p)$ respectively. Write 
        $$\omega_{t_1}=e_{s_1}e_{s_1+1}\cdots e_{k_1} \textit{ \ and \ } \omega_{t_2}=e_{s_2}e_{s_2+1}\cdots e_{k_2}$$
        for some $s_1\leq j_1\leq k_1$ and $s_2\leq j_2\leq k_2$.        Take $$v_1=e_{s_1}e_{s_1+1}\cdots e_{j_1} \text{ and }v_2=e_{s_2}e_{s_2+1}\cdots e_{j_2}.$$ Then we have 
        \begin{enumerate}
            \item $|v_1|\leq|w_{t_1}|\leq d\text{ and }|v_2|\leq|w_{t_2}|\leq d$;
            \item $\omega_{[1,t_1)}v_1\equiv\omega_1\cdots \omega_{t_1-1}e_{s_1}\cdots e_{j_1}\equiv e_1\cdots e_{s_1-1} \cdot e_{s_1}\cdots e_{j_1}\equiv \xi$;
            \item $
        v_1^{-1}\omega_{[t_1,t_2)}v_2\equiv (e_{s_1}\cdots e_{j_1})^{-1}\omega_{t_1}\cdots \omega_{t_2-1}(e_{s_2}\cdots e_{j_2})\equiv b;$
        \item $v_2^{-1}\omega_{[t_2,p]}\equiv (e_{s_2}\cdots e_{j_2})^{-1}\omega_{t_2}\cdots \omega_p\equiv e_{j_2+1}\cdots e_l\equiv \xi^{-1}$.
        \end{enumerate}
        It follows that $1\leq t_1\leq t_2\leq p$ and $\xi,\ v_1,\ v_2\in\textnormal{PSL}(2,\mathbb{Z})$ and $h=b$ are desired. 
        
        If $f_t=c$ or $c^2$, take $\xi=\xi_0$. Then 
        $$\omega\equiv f_1\cdots f_tbf_t^{-1}\cdots f_1^{-1},$$
        where the word in the right hand side is reduced. With similar argument to above, one may find desired $1\leq t_1\leq t_2\leq p$ and $\xi,\ v_1,\ v_2\in\textnormal{PSL}(2,\mathbb{Z})$ and $h\in\{b,c,c^2\}$.

        The proof of Part $(1)$ is complete.
        
        (2). By assumption there exists $\xi_0\in\textnormal{PSL}(2,\mathbb{Z})$ such that $$\omega\equiv \xi_0 h_0^m\xi_0^{-1}$$ where
        $$h_0=f_1\cdots f_t$$
        is a $(b,c)-$representation of $h_0$, in which each $f_s\in\{b,c,c^{2}\}$ and $f_1=b$, $f_t=c$ or $c^{2}$. Take a $(b,c)-$representation of $\xi_0$ whose last letter could be one of $\{b,c,c^2\}$. We only consider the case that the last letter is $b$; the proofs for the cases of $c$ and $c^2$ are similar, that we leave to interested readers. Recall that $b^2 \equiv 1$. Since $f_1=b$ and the last letter of $\xi_0$ is $b$, there are certain cancellations in $\xi_0 h_0^m\xi_0^{-1}.$ So up to maximal cancellations, $\xi_0$ has a $(b,c)-$representation with one of the following forms,
        \begin{enumerate}
            \item $\xi_0=\left(f_t^{-1}\cdots f_1^{-1}\right)^q$ for some $q\geq 0$;
            \item $\xi_0=f_j^{-1}\cdots f_1^{-1}\left(f_t^{-1}\cdots f_1^{-1}\right)^q$ for some $1\leq j\leq t-1$ and $q\geq 0$;
            \item $\xi_0=g_1\cdots g_if_j^{-1}\cdots f_1^{-1}\left(f_t^{-1}\cdots f_1^{-1}\right)^q$ for some $i\geq 1$, $1\leq j\leq t-1$ and $q\geq 0$; moreover, $g_i \neq f_{j+1}^{-1}$. In this case, it is not hard to see that $g_i=f_{j+1}=c$ or $c^2$.
        \end{enumerate}
         We prove the lemma when $\xi_0$ is of form (c) above, i.e.,  
         $$\xi_0=g_1\cdots g_if_j^{-1}\cdots f_1^{-1}\left(f_t^{-1}\cdots f_1^{-1}\right)^q;$$
         the proofs are similar for the other two cases. Then we have
        \begin{align}\label{e-word}
        e_1e_2\cdots e_l=\omega \equiv g_1\cdots g_i f_{j+1}\cdots f_t\left(f_1\cdots f_t\right)^{m-1}f_1\cdots f_j g_i^{-1}\cdots g_1^{-1}.
        \end{align}
        Take $\xi=g_1\cdots g_i$ and $h=f_{j+1}\cdots f_t f_1\cdots f_j$. So we have 
        \[e_1e_2\cdots e_l=\omega \equiv \xi h^m \xi^{-1}.\]
        \begin{itemize}
        \item If $g_i=f_{j+1}=c$, then the word in the right hand side of \eqref{e-word} is reduced. It follows that there exist $1\leq j_1\leq j_2\leq j_3\leq j_4\leq l$ such that
        $$\begin{matrix}e_1\cdots e_{j_1}\equiv \xi, & e_1\cdots e_{j_2}\equiv \xi h,
        \\
        e_1\cdots e_{j_3}\equiv \xi h^2, & e_1\cdots e_{j_4}\equiv \xi h^m.\end{matrix}$$
        Then the claim follows from a similar argument to that of Part $(1)$.

        \item If $g_i=f_{j+1}=c^2$, then $\omega$ could be written as a reduced word of form 
        $$\omega\equiv g_1\cdots g_{i-1}cf_{j+2}\cdots f_t\left(f_1\cdots f_t\right)^{m-1}f_1\cdots f_j g_i^{-1}\cdots g_1^{-1}.$$
        Note that
        \[ g_1\cdots g_{i-1}cf_{j+2}\cdots f_t \cdot f_1\cdots f_j \equiv \xi h.\]
        It follows that there exist $1\leq j_1\leq j_2\leq j_3\leq j_4\leq l$ such that
        $$\begin{matrix}e_1\cdots e_{j_1}c\equiv g_1\cdots g_{i-1}c^2= \xi, & \ e_1\cdots e_{j_2}\equiv \xi h,
        \\
        e_1\cdots e_{j_3}\equiv \xi h^2, & e_1\cdots e_{j_4}\equiv \xi h^m.\end{matrix}$$
        Then the claim follows from a similar argument to that of Part $(1)$.
        \end{itemize}
        
        The proof of Part $(2)$ is also complete.
    \end{proof}

\end{document}
